\title{RWKV: Reinventing RNNs for the Transformer Era}

\begin{abstract}
\label{abstract}
Transformers have revolutionized almost all natural language processing (NLP) tasks but suffer from memory and computational complexity that scales quadratically with sequence length. In contrast, recurrent neural networks (RNNs) exhibit linear scaling in memory and computational requirements but struggle to match the same performance as Transformers due to limitations in parallelization and scalability. We propose a novel model architecture, Receptance Weighted Key Value (RWKV), that combines the efficient parallelizable training of transformers with the efficient inference of RNNs.

Our approach leverages a linear attention mechanism and allows us to formulate the model as either a Transformer or an RNN, thus parallelizing computations during training and maintains constant computational and memory complexity during inference. We scale our models as large as 14 billion parameters, by far the largest dense RNN ever trained, and find RWKV performs on par with similarly sized Transformers, suggesting future work can leverage this architecture to create more efficient models. This work presents a significant step towards reconciling trade-offs between computational efficiency and model performance in sequence processing tasks.
\footnote{Code at: \href{https://github.com/BlinkDL/RWKV-LM}{https://github.com/BlinkDL/RWKV-LM}}
\end{abstract}

\section{Introduction} \label{introduction}
Deep learning has greatly advanced artificial intelligence, impacting a range of scientific and industrial uses. These often involve complex sequential data processing tasks such as natural language understanding, conversational AI, time-series analysis, and indirectly sequential formats like images and graphs \cite{brown2020language, ismail2019deep, wu2020comprehensive,albalak-etal-2022-feta}. Predominant among these techniques include RNNs and Transformers \cite{vaswani2017attention}, each with specific benefits and drawbacks. RNNs require less memory, particularly for handling long sequences.  However, they suffer from the vanishing gradient problem and non-parallelizability in the time dimension during training, limiting their scalability \cite{hochreiter1998vanishing,le2016quantifying}. 

Transformers emerged as a powerful alternative, adept at managing local and long-range dependencies and supporting parallelized training \cite{tay2022efficient}. Models such as GPT-3 \cite{brown2020language}, ChatGPT \cite{openai2022chatgpt,kocon2023chatgpt}, LLaMA \cite{touvron2023llama}, and Chinchilla \cite{hoffmann2022training} showcase the potential of Transformers in NLP. However, the self-attention mechanism's  quadratic complexity makes it computationally and memory intensive for tasks involving long sequences and constrained resources. This has stimulated research to enhance Transformers' scalability, sometimes sacrificing some of their effectiveness~\cite{wang2020linformer,zaheer2020bigbird,dao2022flashattention}. 

To tackle these challenges, we introduce the Receptance Weighted Key Value (\textbf{RWKV}) model, combining the strengths of RNNs and Transformers while circumventing key drawbacks. RWKV alleviates  memory bottleneck and quadratic scaling associated with Transformers  \cite{katharopoulos2020transformers} with efficient linear scaling, while maintaining the expressive properties of the Transformer, such as  parallelized training and robust scalability. RWKV reformulates the attention mechanism with a variant of linear attention, replacing traditional dot-product token interaction with more effective channel-directed attention. This implementation, \emph{without approximation}, offers the lowest computational and memory complexity; see Table~\ref{tab:comparison}.

The motivation behind RWKV is to balance computational efficiency with expressive capacity in neural networks. It offers a solution for handling large-scale models with billions of parameters, exhibiting competitive performance at a reduced computational cost. Experiments suggest RWKV addresses  scaling and deployment challenges in AI, especially for sequential data processing, pointing towards more sustainable and efficient AI models.

Our contributions in this paper are as follows:

\begin{itemize}
 \item The introduction of RWKV, a novel architecture combining RNNs and Transformer advantages while mitigating their limitations.

 \item Detailed experiments demonstrating RWKV's performance and efficiency on benchmark datasets for large-scale models. 
 \item The release of pretrained models, from 169 million to 14 billion parameters, trained on the Pile \cite{gao2020pile,biderman2022datasheet}.\footnote{\href{https://huggingface.co/RWKV}{https://huggingface.co/RWKV}}
\end{itemize}

\section{Background} \label{background}

Here we briefly review the fundamentals of RNNs and Transformers. 

\subsection{Recurrent Neural Networks (RNNs)}

Popular RNN architectures such as LSTM \cite{RNN_LSTM97} and GRU \cite{RNN_GRU2014} are characterized by the following formulation (shown for LSTM, others can be reasoned similarly):

\begin{align}
f_t &= \sigma_g(W_{f} x_t + U_{f} h_{t-1} + b_f), \label{eq:lstm-1}\\
i_t &= \sigma_g(W_{i} x_t + U_{i} h_{t-1} + b_i), \\
o_t &= \sigma_g(W_{o} x_t + U_{o} h_{t-1} + b_o), \\
\tilde{c}_t &= \sigma_c(W_{c} x_t + U_{c} h_{t-1} + b_c), \\
c_t &= f_t \odot c_{t-1} + i_t \odot \tilde{c}_t,\\
h_t &= o_t \odot \sigma_h(c_t). \label{eq:lstm-6}
\end{align}

Although RNNs can be factored into two linear blocks ($W$ and $U$) and an RNN-specific block (\ref{eq:lstm-1})--(\ref{eq:lstm-6}), as noted by \citet{RNN_QuasiRNN17},

the data dependency relying on previous time steps prohibits parallelizing these typical RNNs.

\subsection{Transformers and AFT} \label{sec:background-attn} 
Introduced by \citet{vaswani2017attention}, Transformers are a class of neural networks that have become the dominant architecture for several NLP tasks. Instead of operating on sequences step-by-step like RNNs, Transformers rely on attention mechanisms to capture relationships between all input and all output tokens:

\begin{align}
\mathrm{Attn}(Q, K, V) = \mathrm{softmax}(QK^\top)V,
\end{align}

where the multi-headness and scaling factor $\frac{1}{\sqrt{d_k}}$ is omitted for convenience. The core $QK^\top$ multiplication is an ensemble of pairwise attention scores between each token in a sequence, which can be decomposed as vector operations:

\begin{align}
\mathrm{Attn}(Q, K, V)_t &= \frac{\sum_{i=1}^{T} e^{q^{\top}_t k_i} \odot v_i}{\sum_{i=1}^{T} e^{q^{\top}_t k_i}}.

\end{align}

AFT \cite{zhai2021attention}, alternately formulates

\begin{align}

\mathrm{Attn}^{+}(W,K,V)_t = \frac{ \sum_{i=1}^{t} e^{w_{t,i}+k_i} \odot v_i}{\sum_{i=1}^{t} e^{w_{t,i}+k_i}},
\end{align}

where $\{w_{t,i}\}\in R^{T\times T}$ is the learned pair-wise position biases, and each $w_{t,i}$ is a scalar.

Inspired by AFT, RWKV takes a similar approach. However, for simplicity, it modifies the interaction weights so that it can be transformed into an RNN. Each $w_{t,i}$ in RWKV is a channel-wise time decay vector multiplied by the relative position and traced backward from current time as it decays:

\begin{align}
    w_{t,i} = -(t-i)w,
\end{align}

where $w\in (R_{\geq 0})^{d}$, with $d$ the number of channels. We require $w$ to be non-negative to ensure that $e^{w_{t,i}} \le 1$ and the per-channel weights decay backwards in time.

\section{RWKV} 
\label{rwkv_model}

The RWKV model architecture is defined by four fundamental elements that are intrinsic to the time-mixing and channel-mixing blocks:

\begin{itemize}
    \item $R$: The \textbf{Receptance} vector acts as the receiver of past information.
    \item $W$: The \textbf{Weight} signifies the positional \\ weight decay vector, a trainable parameter within the model.
    \item $K$: The \textbf{Key} vector performs a role analogous to $K$ in traditional attention mechanisms. 
    \item $V$: The \textbf{Value} vector functions similarly to $V$ in conventional attention processes.
\end{itemize}

These core elements interact multiplicatively at each timestep, as depicted in Figure \ref{fig:blocks}.

\subsection{Architecture}

The RWKV model is composed of stacked residual blocks. Each block consists of a time-mixing and a channel-mixing sub-block, embodying recurrent structures to leverage past information.

This model uses a unique attention-like score update process, which includes a time-dependent softmax operation improving numerical stability and mitigating vanishing gradients (for rigorous proof, see Appendix \ref{sec:appendix-gradstab}). It ensures that the gradient is propagated along the most relevant path. Additionally, layer normalization \cite{ba2016layer} incorporated within the architecture aids in stabilizing the gradients, effectively addressing both vanishing and exploding gradient issues. These design elements not only enhance the training dynamics of deep neural networks but also facilitate the stacking of multiple layers, leading to superior performance over conventional RNN models by capturing complex patterns across different levels of abstraction (see also Appendix \ref{sec-appen:vis}).

\subsubsection{Token Shift} \label {tokenshift}

In this architecture, all linear projection vectors ($R$, $K$, $V$ in time-mixing, and $R'$, $K'$ in channel-mixing) involved in computations are produced by linear interpolation between current and previous timestep inputs, facilitating a token shift.

The vectors for time-mixing computation are linear projections of linear combinations of the current and previous inputs of the block:

\begin{align}
\label{eq:time-mix1}
r_t &= W_r \cdot (\mu_r \odot x_t + (1 - \mu_r) \odot x_{t-1} ), \\
\label{eq:time-mix2}
k_t &= W_k \cdot (\mu_k \odot x_t + (1 - \mu_k) \odot x_{t-1} ), \\
\label{eq:time-mix3}
v_t &= W_v \cdot (\mu_v \odot x_t + (1 - \mu_v) \odot x_{t-1} ),
\end{align}

as are the channel-mixing inputs:

\begin{align}\label{eq:channel_mix_eqs1}
r'_t &= W'_r \cdot (\mu'_r \odot x_t + (1 - \mu'_r) \odot x_{t-1} ), \\
\label{eq:channel_mix_eqs2}
k'_t &= W'_k \cdot (\mu'_k \odot x_t + (1 - \mu'_k) \odot x_{t-1} ).
\end{align}

The token shift is implemented as a simple offset in the temporal dimension at each block using the PyTorch \citep{paszke2019pytorch} library as \texttt{nn.ZeroPad2d((0,0,1,-1))}.

\subsubsection{WKV Operator}

The computation of the $WKV$ operator in our model parallels the method used in Attention Free Transformer (AFT) \cite{zhai2021attention}. However, unlike AFT where $W$ is a pairwise matrix, our model treats $W$ as a channel-wise vector that is modified by relative position. In our model, this recurrent behavior is defined by the time-dependent update of the $WKV$ vectors, formalized in the following equation:

\begin{align}
    \label{eq:nom-denom}
    wkv_t &= \frac{ \sum_{i=1}^{t-1} e^{-(t-1-i)w+k_i} \odot v_i + e^{u+k_t} \odot v_t }{\sum_{i=1}^{t-1} e^{-(t-1-i)w+k_i} + e^{u+k_t}}.
\end{align}

To circumvent any potential degradation of $W$, we introduce a vector $U$ that separately attends to the current token. More information about this can be found in Appendix \ref{sec-appen:vis}.

\subsubsection{Output Gating}

Output gating is implemented in both time-mixing and channel-mixing blocks using the sigmoid of the receptance, $\sigma(r)$. The output vector $o_t$ post the $WKV$ operator is given by:

\begin{align}
\label{eq:time-mix4}
o_t &= W_o \cdot (\sigma(r_t) \odot wkv_t).
\end{align}

In the channel-mixing block, a similar operation is performed:

\begin{align}
\label{eq:channel_mix_eqs3}
o'_t &= \sigma(r'_t) \odot (W'_v \cdot \max({k'_t}, 0)^2),
\end{align}

where we adopt the squared ReLU activation function \cite{DBLP:journals/corr/abs-2109-08668}.

\subsection{Transformer-like Training} \label{transformer_parallel}

RWKV can be efficiently parallelized using a technique called \textit{time-parallel mode}, reminiscent of Transformers. The time complexity of processing a batch of sequences in a single layer is $O(BTd^2)$, primarily consisting of matrix multiplications $W_\lambda$, where $\lambda\in \{r,k,v,o\}$ (assuming $B$ sequences, $T$ maximum tokens, and $d$ channels).

In contrast, updating attention scores $wkv_t$ involves a serial scan (see Appendix \ref{sec:appendix-rnn} for more detail) and has complexity $O(BTd)$.

The matrix multiplications can be parallelized similarly to $W_{\lambda}$, where $\lambda \in \{Q,K,V,O\}$ in conventional Transformers. The element-wise $WKV$ computation is time-dependent but can be readily parallelized along the other two dimensions \cite{lei-etal-2018-simple}\footnote{For extremely long sequences, more sophisticated methods such as \citet{Martin2017ParallelizingLR} that parallelize over sequence length could be used.}. 

\subsection{RNN-like Inference} \label{rnn_inference}

Recurrent networks commonly utilize the output at state $t$ as input at state $t+1$. This usage is also observed in the autoregressive decoding inference of language models, where each token must be computed before being passed to the next step. RWKV takes advantage of this RNN-like structure, known as \textit{time-sequential mode}. In this context, RWKV can be conveniently formulated recursively for decoding during inference, as demonstrated in Appendix~\ref{sec:appendix-rnn}.

\subsection{Additional Optimizations} \label{adding_opts}

\paragraph{Custom Kernels} \label{custom_kernel}

To address inefficiencies in the $WKV$ computation arising from the sequential nature of the task when using standard deep learning frameworks, we have developed a custom CUDA kernel. This kernel enables the execution of a single compute kernel on training accelerators, while all other parts of the model, such as matrix multiplications and point-wise operations, are already inherently parallelizable and efficient. 

\paragraph{Small Init Embedding} \label{smallinit}

During the initial stage of training a transformer model \cite{vaswani2017attention}, we observe that the embedding matrix undergoes slow changes, presenting a challenge for the model to move away from its initial noisy embedding state. To address this issue, we propose an approach that involves initializing the embedding matrix with small values and subsequently applying an additional LayerNorm operation. This accelerates and stabilizes the training process, allowing for the training of deep architectures with post-LN components. The effectiveness of this approach is demonstrated in Figure \ref{fig:small_init_emb}, illustrating improved convergence by enabling the model to quickly transition away from the initially small embedding. This is achieved through small changes occurring in a single step, which subsequently lead to substantial alterations in directions and further notable changes after the LayerNorm operation.

\paragraph{Custom Initialization}

Building on principles from previous works \citep{he2016identity, af2_2021_nature}, we adopt an initialization strategy where parameters are set to values resembling an identity mapping while breaking symmetry to establish a clear information flow. The majority of weights are initialized to zero, and linear layers do not employ biases. Detailed formulas are given in Appendix \ref{sec:initialization_formulas}. We observe that the choice of initialization plays a crucial role in both the speed and quality of convergence (refer to Appendix \ref{appendix_smallinit} for further details).

\subsection{Implementation}\label{detailed_implementation}

RWKV is implemented using the PyTorch Deep Learning Library \cite{paszke2019pytorch}. We integrate additional optimization strategies inspired by DeepSpeed \cite{10.1145/3394486.3406703} into the system, improving its efficiency and scalability.

The model begins with an embedding layer, as detailed in Section \ref{smallinit}. Following this are several identical residual blocks arranged sequentially. These are depicted in Figures \ref{fig:blocks} and \ref{fig:arch} and  adheres to the principles outlined in Section \ref{tokenshift}. After the last block, a simple output projection head, consisting of a LayerNorm \cite{ba2016layer} and a linear projection, is employed for logits generation for next-token prediction and  computation of the cross-entropy loss during training.

\section{Trained Models and Computing Costs}\label{appendix:flop_count}

To demonstrate the scalability of RWKV, we train six models ranging from 169 million to 14 billion parameters as shown in Table \ref{tab:model_flop_count}. All models are trained for one epoch (330 billion tokens) on the Pile \citep{gao2020pile,biderman2022datasheet}.

The number of parameters for each model is computed using the formula:  $\text{\# parameters} = 2VD + 13D^2L + D(11L+4)$ where $V$ = 50277 is the vocabulary size, $D$ represents the Model Dimension and $L$ corresponds to the number of layers. FLOPs is for a forward pass for one token. It was calculated as $2(2VD + 13D^2L)$, which is the twice (add and multiply) the number of parameters in linear layers. The backwards pass FLOPs can be approximated as twice that of the forward pass, giving a total of $6(2VD + 13D^2L)$ FLOP per token. Notably, this matches the standard formula for FLOP calculations in transformers \citet{kaplan2020scaling}: $\text{FLOP} = 6\cdot [\text{\# tokens}]\cdot [\text{\# parameters}]$.

\subsection{Additional Training Details}

For training, we use the standard Adam optimizer without weight decay, use bfloat16 precision, and train with a context length of 1024 tokens. Further details on hyperparameters are in Appendix~\ref{appendix_hyperparam}. Diverting from standard practice for transformers, we apply exponential decay to our learning rate. We also incorporate the auxiliary loss introduced by PaLM \cite{PaLM}, supplementing the standard cross-entropy loss function. This auxiliary loss encourages the softmax normalizer to approximate zero closely. As for the learning rate schedule, it remains constant for the initial iterations, and subsequently decays exponentially. 

\subsection{Scaling Laws}

Scaling laws \cite{kaplan2020scaling,henighan2020scaling,hoffmann2022training,muennighoff2023scaling} in language models refer to the mathematical relationships that describe how the performance of a language model changes with respect to various factors. These factors can include the model size ($N$), dataset size ($D$), or the optimally allocated compute budget ($C_{\rm min}$). Scaling laws are important for two primary reasons: they allow us to make predictions and plans regarding the costs and performance of large models before they are trained via interpolation and extrapolation \cite{black2022gpta,le2022language} and the contexts in which they fail provides rich feedback on important areas for future research \cite{Wei2022EmergentAO,biderman2023emergent}.

Previous work on scaling laws for RNNs has claimed that LSTMs do not strictly follow the same log-log linear scaling that transformers do \cite{kaplan2020scaling}. We train 45 RWKV models for a variety of pairs (dataset, parameters) and find that RWKV \textit{does} follow the same general form of the scaling law that is well established for transformers. Figure \ref{fig:scaling} shows our results for loss as a function of compute, with the linear fit to the Pareto optimal points holding an $r^2$ value of $0.994$. Even when we extrapolate our curve an additional order of magnitude (blue), we find an extremely good fit with an $r^2$ of $0.875$.

\section{Evaluations}\label{evaluations}

Having demonstrated the scalability of RWKV models in the previous section, we now turn our attention to their competitiveness with traditional transformers. We focus on two questions:

\paragraph{Competitiveness} Is RWKV competitive against quadratic transformer architectures with the same amount of compute?

\paragraph{Long Context} Does increasing the context length of RWKV yield better language modeling loss when RWKV models are trained for context lengths that most open-sourced quadratic transformers \textit{cannot} efficiently process?

\subsection{NLP Evaluations}

To demonstrate that RWKV is competitive with traditional transformers at NLP tasks, we compare with similarly sized models trained for a similar number of tokens (Pythia \cite{biderman2023pythia}, OPT \cite{zhang2022opt} and BLOOM \cite{scao2022bloom}). All RWKV models were trained for one epoch on the Pile (330B tokens), which is close but not identical to the amount of tokens the Pythia, OPT, and BLOOM models were trained for. Consequently, we compare our models on a \textit{FLOP-matched basis}. We avoid comparing with model trained in the Chinchilla-optimal regime \citep{hoffmann2022training} or the overtrained regime \citep{touvron2023llama} to ensure the most equitable comparison.

We report results on ARC (both Easy and Challenge) ~\cite{ARC2018}, BoolQ ~\cite{clark2019boolq}, COPA ~\cite{COPA2011}, HeadQA ~\cite{HeadQA2020}, HellaSwag ~\cite{HellaSwag2019}, LAMBADA~\cite{LAMBADAdataset}, OpenBookQA ~\cite{OpenBookQA2018}, PIQA~\cite{Bisk2020}, ReCoRD ~\cite{ReCord}, SciQ ~\cite{SciQ2017}, and Winogrande ~\cite{Wino2020}. Figure \ref{fig:average} shows the average results across all benchmarks. Some individual benchmarks are shown in Fig \ref{fig:eval}, with the rest in Appendix \ref{sec-appen:eval-detail}.

Additionally, we carried out comparative studies on RWKV and ChatGPT / GPT-4, see Appendix \ref{appendix:promptImportanceAndRWKVvsGPT}. They revealed that RWKV is very sensitive to prompt engineering. When the prompts were adjusted (re-ordered) from the ones used for GPT to more suitable for RWKV, the performance (F1) increased even from 44.2\% to 74.8\%. For sarcasm detection, RWKV outperformed ChatGPT, but was still slightly worse than the SOTA solution.

\subsection{Extended Context Finetuning}

Unlike transformers, RNNs do not have a pre-defined sequences length when they are created. However in order to efficient make use of compute we nevertheless need to preprocess the training data into contexts of the same length. We find that we are able to teach the model how to efficiently handle substantially larger batch sizes by finetuning with progressively increasing sequence length. Specifically, we first double the sequence length from 1024 to 2048 and finetune for 10B tokens from the original pretraining corpus, then we double again to 4096 for 100B tokens from the same corpus, and finally double to 8192 tokens for another 100B tokens from the same corpus. In  Fig.~\ref{fig:ctxlen_rwkv_loss} we show that increasing context length leads to lower test loss on the Pile, an indication that RWKV can make effective use of long contextual information.

\subsection{Long Context Benchmarks}

Additionally, we evaluate our model's ability to handle very long sequences by comparing to state-of-the-art long sequence models on the Long-Range Arena (LRA) benchmark \cite{tay2021long}. LRA is designed to assess the performance of models in handling lengthy context situations. It includes a collection of tasks with sequences ranging from 1,000 to 16,000 tokens, covering various types of data like text, natural language, synthetic images, and mathematical expressions. We apply RWKV on the LRA benchmark and the results are in Appendix \ref{lorfa}. The results show that RWKV performs second only to the S4 model in five datasets.

\begin{comment}
\section{Scaling Laws} \label{scaling laws}

Scaling laws \cite{kaplan2020scaling,henighan2020scaling,hoffmann2022training,muennighoff2023scaling} in language models refer to the mathematical relationships that describe how the performance of a language model changes with respect to various factors. These factors can include the model size ($N$), dataset size ($D$), or the optimally allocated compute budget ($C_{\rm min}$). Scaling laws are important for two primary reasons

\begin{enumerate}
    \item They allow us to make predictions and plans regarding the costs and performance of large models before they are trained via interpolation and extrapolation \cite{black2022gpta,le2022language}
    \item The contexts in which they fail provides rich feedback on important areas for future research \cite{Wei2022EmergentAO,biderman2023emergent}
\end{enumerate}

We have observed that the RWKV model also follows scaling laws and the test loss $L$ can be estimated using a power-law relationship. This observation is depicted in Figure \ref{fig:scaling-laws}.

We conducted experiments on a range of RWKV models with different model sizes, ranging from 169M to 14B parameters. In doing so, we discovered the first scaling law:
\begin{align}
L(N) = \left(N_{\mathrm{c}}/N\right)^{\alpha_N}, \label{eq:scale_law_N}
\end{align}
where $\alpha_N \sim 0.072$ and $N_{\mathrm{c}} \sim 3.1 \times 10^{13}$ (non-embedding parameters). This scaling law demonstrates that as the model size increases, the test loss decreases without encountering any bottlenecks or plateaus, indicating that the RWKV architecture exhibits excellent scalability.

When the dataset size increases by an order of magnitude, we also observe a linear decrease in loss, which represents the second scaling law:
\begin{align}
L(D) = \left(D_{\mathrm{c}}/D\right)^{\alpha_D}, \label{eq:scale_law_N_2}
\end{align}
where $\alpha_D \sim 0.066$ and $D_{\mathrm{c}} \sim 1.4 \times 10^{15}$ (tokens).

The amount of computation is determined by two factors: the dataset size and the model size. Therefore, when there is an ample amount of data and an appropriate model size, the computational cost also follows the scaling law:
\begin{align}
L(C) = \left(C_{\mathrm{c}}/C\right)^{\alpha_C} \label{eq:scale_law_N_3},
\end{align}
where $\alpha_C \sim 0.053$ and $C_{\mathrm{c}} \sim 1.1 \times 10^{7}$ (PF-days). This indicates that with sufficient computational resources, combined with an ample amount of data and an appropriate model size, the RWKV model can continue to scale.

The next question we attempt to answer is how to make the trade-off between model size and dataset size when only a limited computing budget is available.

\end{comment}

\section{Inference Experiments} \label{infer_experiments}

We benchmark inference requirements according to size and family. Specifically, we evaluate text generation speed and memory requirements on typical compute platforms including CPU (x86) and GPU (NVIDIA A100 80\,GB). For all of our inference experiments we use float32 precision and the HuggingFace Transformers~\cite{hf_transformers}. We include all model parameters in the parameter count, including both embedding and non-embedding layers. Performance under different quantization setups is left to further work. See Appendix \ref{appendix:inference_results} for more results.

\section{Future Work} \label{future_work}

There are several promising directions for future work on the RWKV architecture. Work can be done to increase model expressivity by enhancing the time-decay formulations and exploring initial model states while maintaining efficiency.

The RWKV computational efficiency can be further improved by applying a parallel scan in the $wkv_t$ step to reduce the computational cost to $O(B\log(T)d)$.

The mechanisms used in RWKV can be applied to encoder-decoder architectures, potentially replacing the cross-attention mechanism. This could be applicable in seq2seq or multimodal settings, thereby enhancing efficiency during both training and inference. 

RWKV's state (or \textit{context}) can be leveraged for interpretability, predictability in sequence data, and safety. Manipulating the hidden state could also guide behavior and allow greater customizability through prompt tuning.

The RWKV architecture is not perfect, and can be improved via many aspects, such as modifying the formulae or implementing larger internal states. Larger states can enhance the model's memory to previous context and improve performance over various tasks.

\section{Conclusions} \label{conclusions}

We introduced RWKV, a new approach to RNN models exploiting the potential of time-based mixing components. RWKV introduces several key strategies that allow it to capture locality and long-range dependencies while addressing limitations of current architectures by: (1) replacing the quadratic QK attention with a scalar formulation at linear cost, (2) reformulating recurrence and sequential inductive biases to enable efficient training parallelization and efficient inference, and (3) enhancing training dynamics using custom initializations. 

We benchmark the proposed architecture in a wide variety of NLP tasks and show comparable performance to SoTA with reduced cost. Further experiments on expressivity, interpretability, and scaling showcase the model capabilities and draw parallels in behavior between RWKV and other LLMs. 

RWKV opens a new route for scalable and efficient architectures to model complex relationships in sequential data. While many alternatives to Transformers have been proposed with similar claims, ours is the first to back up those claims with pretrained models with tens of billions of parameters.

\section{Limitations}\label{limitations}
While our proposed RWKV model has demonstrated promising results regarding training and memory efficiency during inference, some limitations should be acknowledged and addressed in future work.

First, the linear attention of RWKV leads to significant efficiency gains but still, it may also limit the model's performance on tasks that require recalling minutiae information over very long contexts. This is due to the funneling of information through a single vector representation over many time steps, compared with the full information maintained by the quadratic attention of standard Transformers. In other words, the model's recurrent architecture inherently limits its ability to ``look back'' at previous tokens, as opposed to traditional self-attention mechanisms. While learned time decay helps prevent the loss of information, it is mechanistically limited compared to full self-attention. 

Another limitation of this work is the increased importance of prompt engineering in comparison to standard Transformer models. The linear attention mechanism used in RWKV limits the information from the prompt that will be carried over to the model's continuation. As a result, carefully designed prompts may be even more crucial for the model to perform well on tasks.

The above RWKV property was confirmed by studies on prompt engineering presented in Appendix \ref{appendix:promptImportanceAndRWKVvsGPT}. By changing the order of the information pieces, we were even able to almost double the RWKV performance for some tasks. 

\section{Ethics Statement}

In this paper, we present a novel architecture for sequential data processing and prove its effectiveness by building a series of LLMs trained on publicly released pretraining data \cite{gao2020pile,biderman2022datasheet} and later fine-tuned on publicly available instructions \cite{alpaca, codealpaca, guanaco, gpt4all, sharegpt, Firefly, BELLE,belle2023exploring}.

As a novel architecture for sequential data, RWKV has the potential to improve sequence-based models across different applications ranging from natural language processing to biomedical data processing or climate modelling. Since the training code is released open source, RWKV contributes to the democratization of AI, levels the playing field, and empowers members of the Open Source community to inspect, study, and finetune RWKV in particular tasks. Moreover, it contributes to advancing the understanding of LLMs capabilities and limitations. A significant amount of work has been devoted to increasing the efficiency of RWKV training  so as to minimize its cost and promote accessibility.

As LLMs trained on public data, RWKV's lower inference cost compared to Transformer alternatives makes it more suitable for deployment in consumer and edge hardware, which is a step towards the democratization and distribution of LLMs to the general public, creating better privacy and ownership incentives. It also lowers the resource barrier to Chat assistants and text generation for small and/or underrepresented communities. PreTrained model weights for different sizes ranging from 0.1B to 14B parameters trained on multiple languages are released to increase ease of adoption and allow for the study of emergent phenomena. 

On the other hand, with lower resource barriers, the spreading of AI-generated text might become more prevalent. Current RWKV LLMs may exhibit and/or reproduce biases and potentially harmful content present in the data used for training. Nonetheless, mitigation and finetuning strategies discussed for other, large Transformer models should be applicable to RWKV as well. 

\end{document}
